% ---
% Inicia os anexos
% ---
\begin{anexosenv}
    \chapter{Código Exercício 5.7}\label{cap:ex5-7-codigo}
    
    \begin{lstlisting}[language=Python, breaklines=true, basicstyle=\tiny]
    import numpy as np
    import matplotlib.pyplot as plt
    import seaborn as sns
    from scipy.fft import fft, fftfreq

    # --- parâmetros físicos ---
    eps = 8.854187817e-12  # Vacuum permittivity (F/m)
    mu = 4 * np.pi * 1e-7  # Vacuum permeability (H/m)
    vp = 1 / np.sqrt(eps * mu)
    f = 1e9

    # --- domínio ---
    min_lambda = vp / f  # comprimento de onda mínimo para 1 GHz = 3e-1 m
    Nx = 100
    Ny = 50
    dx = dy = 0.1 * min_lambda
    dt_list = [0.5, 0.9, 0.98, 0.99, 1, 1.01, 1.1] / (vp * np.sqrt((1/dx**2) + (1/dy**2)))
    # dt_list = [0.99] / (vp * np.sqrt((1/dx**2) + (1/dy**2)))
    max_value_per_dt = []

    for dt in dt_list:
        Nt = 500
        L = Nx * dx
        T = Nt * dt

        grid_x_Ez = np.linspace(0, L, Nx + 1)
        grid_y_Ez = np.linspace(0, L, Ny + 1)
        grid_t_Ez = np.linspace(0, T, Nt + 1)

        grid_x_Hx = np.linspace(0, L, Nx + 1)
        grid_y_Hx = np.linspace(dy/2, L - dy/2, Ny)
        grid_t_Hx = np.linspace(dt/2, T - dt/2, Nt)

        grid_x_Hy = np.linspace(dx/2, L - dx/2, Nx)
        grid_y_Hy = np.linspace(0, L, Ny + 1)
        grid_t_Hy = np.linspace(dt/2, T - dt/2, Nt)

        # --- campos ---
        Ez_curr = Ez_next = np.zeros((grid_x_Ez.shape[0], grid_y_Ez.shape[0]))
        Hx_curr = Hx_next = np.zeros((grid_x_Hx.shape[0], grid_y_Hx.shape[0]))
        Hy_curr = Hy_next = np.zeros((grid_x_Hy.shape[0], grid_y_Hy.shape[0]))

        # valores do tempo para salvar os snapshots e montar os heatmaps
        n_snapshots = np.linspace(5, Nt - 5, 6, dtype=int)
        Ez_snapshots = []

        # ponta de prova no centro
        probe = []
        cx, cy = Nx // 2, Ny // 2

        for n in range(Nt-1):
            # fonte onda planar
            source = np.exp(-((n * dt - 1.8e-9)**2) / (0.6e-9)**2)
            source_senoidal = np.sin(2 * np.pi * f * n * dt)

            # aplica em todos os Y em X = 0 (left boundary plane wave)
            Ez_curr[0, :] = source_senoidal

            Hx_next[:, :] = Hx_curr[:, :] - (dt/(mu*dy)) * (
                Ez_curr[:, 1:] - Ez_curr[:, :-1]
            )

            Hy_next[:, :] = Hy_curr[:, :] + (dt/(mu*dx)) * (
                Ez_curr[1:, :] - Ez_curr[:-1, :]
            )

            Ez_next[1:-1, 1:-1] = Ez_curr[1:-1, 1:-1] + (dt/eps) * (
                (Hy_next[1:, 1:-1] - Hy_next[:-1, 1:-1]) / dx -
                (Hx_next[1:-1, 1:] - Hx_next[1:-1, :-1]) / dy
            )

            # PEC an extremidade direita
            Ez_next[-1, :] = 0

            # periodic boundary condition
            Ez_next[:, 0] = Ez_next[:, Ny]
            Hx_next[:, 0] = Hx_next[:, Ny - 1]

            # atualiza os campos para o próximo passo
            Ez_curr = Ez_next.copy()
            Hx_curr = Hx_next.copy()
            Hy_curr = Hy_next.copy()

            # salva na ponta de prova
            probe.append(Ez_curr[cx, cy])

            # salva para os snapshots
            if n in n_snapshots:
                Ez_snapshots.append(Ez_curr.copy())

        max_value_per_dt.append(np.max(probe))

    \end{lstlisting}

\end{anexosenv}
